\textbf{Figure 5. Liver protein tests the disease-state RNA changes, and intact tissue places Figure 5 genes in cell types.}

\textbf{(A)} Liver protein against bulk mRNA change for 3,054 measured genes of a 10,044-gene prioritized universe (after excluding multi-gene semicolon protein groups, leaving 5,836 uniquely mapped proteins). The universe is partly selected on RNA: rho is 0.487 in RNA-selected genes and 0.305 in the 664 genes selected on genetic evidence alone. Protein model: 46 MASLD versus 12 control participants, adjusted for acquisition batch, age, BMI and sex. Blue, protein-significant (BH p \textless{} 0.05 over 5,836 proteins) with the same sign as mRNA; magenta, opposite sign; gray, not significant. The bar shows the 211 protein-significant genes: 188 with the same sign as mRNA (89.1\%) and 23 with the opposite sign (10.9\%). Outlined, the eight recurring example genes (Figure S2AC) (protein log2FC): protein-significant with the same sign as mRNA, CFHR5 -0.79, ADH1B -0.41, BCKDHB -0.29 and DBT -0.24; protein-significant with the opposite sign, BCKDHA -0.25 (mRNA +0.11, not significant); protein not significant, PPM1K -0.70, CFHR4 -0.33 and HKDC1 +0.50.

\textbf{(B)} 25 proteins fixed before the bulk refit, grouped by process, evaluated on the unified MASLD-only proteomics pipeline (5,836 uniquely mapped proteins; 9 canonical RNA calls, 23 protein calls). The same cohort helped select them, so the view is descriptive. Heatmap, protein z-scores ordered by fibrosis stage then NAS. Middle pie-chart grid, partial Spearman rho with each histology score (steatosis, ballooning, inflammation, fibrosis, NAS) within 46 MASLD participants adjusted for batch, age, BMI and sex; filled sector angle and colour encode signed partial Spearman rho and faded sectors mark BH q $\ge$ 0.05 over 125 tests. Right, filled markers are calls (BH padj \textless{} 0.05, plus absolute log2FC \textgreater{} 0.5 for mRNA). Process adjusted NES and BH q from preranked fgsea.

\textbf{(C)} Cell type that makes each gene in intact healthy liver. Xenium in situ data of Tzouanas et al.~(2026; 477-gene panel), 9 control livers; the 9 MASH-cirrhosis explants of the same source are not used here (spillover-corrected using KNG1 and FBP1 reference genes). Rows: the 14 on-panel genes of panel A with a protein change in the same direction as RNA (lower or higher in MASLD), and the 11 on-panel colocalized genes detected in at least 5\% of their top cell type's cells in at least 7 of 9 controls (SPI1 is not detected). Dots: spillover-corrected counts per cell, averaged over donors with equal weight, scaled to the gene's top cell type. Endothelial, stellate/smooth-muscle, macrophage and T/NK columns show the highest of their subtypes. Right, the gene's top snRNA lineage in the deduplicated atlas; ``same'' when it falls in the top tissue cell type (12 of 14 and 8 of 11). Cell labels and coordinates are the source's. Descriptive; no test.

\textbf{(D)} The 21 of 96 scorable fixed programs whose liver protein score is robustly associated with histologic burden in 46 MASLD participants (unified MASLD-only proteomics pipeline, 5,836 uniquely mapped proteins): partial Spearman rho adjusted for acquisition batch, age, BMI and sex, with tie-free Dirichlet participant-bootstrap 95\% CI. Robust: Freedman-Lane BH q \textless{} 0.05 over the 96 scorable programs and stable sign across score variants and both acquisition batches. Rows are selected on the plotted statistic; all 117 programs are in the source table. Blue, lower with greater burden (13); red, higher (8). Histologic burden is the first principal component of steatosis, ballooning, lobular inflammation and fibrosis stage.
